\documentclass[12pt]{article}

\usepackage[a4paper, total={6in, 8in}]{geometry}
\usepackage{textcomp}

\begin{document}
\setlength{\parindent}{0pt}

\def \t {\textrightarrow}
\def \v {\vspace{1em}}
\def \bi {\begin{itemize}}
\def \ei {\end{itemize}}
\def \s[#1] {\section*{#1}}
\def \ss[#1] {\subsection*{#1}}
\def \sss[#1] {\subsubsection*{#1}}

\s[Montale]
\ss[Forse un mattino andando in un'aria di vetro]
Il testo inizia con un dubbio \\
Uomini che non si voltano \t la comprensione del mistero dell'esistenza, e l'averlo sentito dentro \t in fondo non è altro che una maledizione \\
Gli uomini che non si pongono domande, che non si guardano dentro \t il non porsi delle domande \\
Capisce che l uomo è solo, che l esistenza è un inganno \t ma questo lo rende solo, perche gli altri non vogliono porsi domande \\
Immagine fa capire il dischiudersi della verita \t e il suo volgersi indietro e il compiersi di un miracolo (miror, miraris = qualcosa di ignoto che soprende) \\
Il miracolo è capire che c'è il nulla, il vuoto \t paura di ubriaco che porta a perde equilibrio \\
Pier Vincenzo Mengaldo \t poesia del 900 \t si coglie la volonta di descrivere lo stato di confusione, e lo scoraggiamento \\
Le immagini dello schermo sono quasi riflesse nella natura, in "Non recidere forbice" \t e anche in "Cigola la carrucola nel pozzo" \\
La nebbia di sempre, si sfolla \t improvvisamente come una carrellata, case e colli \t la natura e il paesaggio vengono accumulati senza virgole (al verso 6) \\
Franco Cortini, "gli alberi non sono alberi" \\
Con "ma" cambia registro, dopo che ha parlato di inganno \t l uomo è abbituato a lasciar scorrere le immegini senza soffermarsi \\
Montale ha capito che la vita è un inganno \t e infatti se ne va zitto, non puo rivelarlo \t il segreto della vita, del nulla della vita \\
Lo svelarsi del miracolo si corella all'"epifany" di Joyce, ovvero la rivelazione, manifestazione \\
Brevitas che caratterizza la lirica \t ossi di seppia 24-26, occasioni del 39 (con focus su irma brandeis con clizia)

\ss[Meriggiare]
Questo è il canto del limite \\
Oggettistica montaliana, qua ha rilievo fondamentale \t il mondo di montale percepito come limitazione all esistere, e impossibilita dell uomo \\
Natura secca, e nonostante le scaglie di mare che palpitano, c'è solo la certezza di un percorso travagliato, e l impossibilita di ? \\
Il male di vivere, qua ha oggettistica ancora + presente \t meriggiare, fa del nome un verbo \\
In elliot si legge l'afa, e sofferenza che si traduce in una sonnolenza \t il meriggiare (all infinito) evoca una atemporalita di questo spettacolo della natura \\
Il meriggiare lascia percepire la calura, si è presso un muro rovente di orto \t orto come limite di sofferenza, hortus conclusus \\
Sono presenti infiniti e gerundi \t che sottolineano la atemporalita e dell immutabilita del dolore umano \\
Ascoltare tra i pruni e gli sterpi \t evoca natura ostile all uomo, non accudente ma limitante \\
Schiocchi, frusci \t onomatopee \t suoni sono ostici quanto il paesaggio, e sono rappresentazione del paesaggio \\
Attenzione alle cose quotidiane \\
Crepe del suolo \t nella desolazione della natura si trova la fragilita dell uomo \t non ha appoggi \\
La formica rappresenta l operosita \t continuano il loro lavoro, incuranti \\
Vedere a stralci \t mare, di leopardiana memoria, che pero qua viene colto nella sua limitatezza e non con la sua immensita \t pascoli, anche il cielo si vede sono attraverso trame \\
Scaglie di mare \t metafora e sinestresia \\
Umilta di un paesaggio ligure \t che vengono presentati come una provocazione a quella stasi dei primi versi \t schricchiolio è onomatopea \\
Mondo quotidiano, che sottointende un dolore universale \t il sole non porta conforto, ma abbaglia \t il sentire dell animo, stesse dinamiche degli infiniti e dei gerundi \\
Emblematicita porta all ossimoro di "triste meraviglia" \\
Il muro d'orto è diventato una muraglia, con in cima i cozzi aguzzi di bottiglia \t cozzi aguzzi denotano l'ostilita della vita, e l immposibilita di andare al di la e di trovare risoluzone del male di vivere, descritto nel paesaggio \\
L'autore ha gia percorso qui un percorso di illuminazione \t ha percepito sofferenza

\ss[Ho sceso dandoti il braccio]
"Satur" \t vuol dire pieno, e satira deriva da questa parola \t la satira era infatti un piatto misto \\
La lirica è indirizzata alla moglie (non a Irma) \t esce nel 72, anno dopo aver sposato la moglie (che muore un anno dopo) \\
Quotidianita apparentemente semplice, trasformata in una complessita \\
La moglie è ipovedente, viene paragonata a una mosca (in altra lirica) \\
Due valori: valore transitivo e intransitivo (di scendere) \t si trova un valore transitivo \\
Un milione di scale è iperbole \t evoca il loro lungo trascorso insieme \\
Qua l'assenza di certezza è ogni gradino \\
L'avrebbe voluta avere ancora \t la mancanza è forte, è breve il nostro lungo viaggio \t il mio dura tutto ora \\
Punto focale: chi crede che la realta sia quella che si vede \t mondo del 900 condensato qua \t Sartre, Freud, Cosini, Pirandello \t qua c'è la fulmineità \\
Questa è la lirica che parte dal personale, ed eleva all'universale \t qua c e descrizione nitida del 900 \\
Prima strofa è + lunga \t tipico del 900, rinnovata armonia tra parola scritta e spazi bianchi \t che si trova anche nel futurismo, ungaretti, etc. \\
Quattr'occhi \t perche erano i suoi di lui e quelli della moglie \t poliptoto, che rafforza questa vita come viaggio \t di noi due, le sole vere pupille (metonimia) \\
Offuscate perche l immagine le appariva tremula, non definita \t la sua inibizoine della vista \\
Testo tardo, tratto da satura \t il sentre del vuoto del 900, che ic arriva mediato attraverso un esperienza soggettiva personale 

\ss[La casa dei doganieri]
Lettura delle emozioni che sono cristallizzate dentro, e che palpitano negli oggetti \\
Si rivolge direttamente a una donna, che non è ne irma ne la moglie \t forse Annetta o Arletta, con cui ha condiviso estati \\
Testo abbastanza prolisso rispetto allo stile montaliano \t e testo dice molto in merito alla sua sensibilita, al suo ricordarsi di qualcuno \\
Usa tono colloquiale \\
La casa si trova su una scogliera, e si trova in liguria \t la casa attende un ritorno che non ci sara \t tema del nostos \\
Sciame, una folla di pensieri \t libeccio è un vento del sud \t che sferza da anni le vecchie mura, con iperbato \\
Sembra di leggere gli occhi che sono spenti, che ricordano quelli di laura in "erano i capei doro a laura sparsi" di Petrarca \\
La bussola è impazzita \t il calcolo dei dadi piu non torna \t descrivono una sorta di ritono dei matematici aristotelici, che anticipa il calcolo combinatorio di Calvino \t Averroe \\
Il calcolo non da certezza \t come non chiederci la formula in "Non chiederci la parola" ? \\
Il filo che tiene insieme i nostri ricordi, si intrica \t io ne tengo ancora un capo, voglio ancora trattenerti \\
Affumicata perche è arsa dalla salsedine e dal sole \t anche lei è priva di vita, perche il sale e sole privano di acqua (vita) \\
La luce di una petroriera \t tema del progresso, letto tramite le navi da carico \t questa luce intermittente da possiblita di speranza? \t il varco è qui, varco è parola focale \\
Riprende con anafora evidente, tu non ricordi \\
La capacita è anche nel percepire il sentimento \t lei pone una distanza, che si coglie anche in cigola la carrucola \\
Mia sera \t enj. esprime un io diviso \t e io non ho piu la dimensione di questo luogo dell anima, e di quella bussola che va all impazzita produce qui questo effetto \t non sa chi rimane e chi va nel suo cuore \\
Conclusione emblematica, che sottointende una delusione \t l'abbandono di lei, mentre lui ha ancora sentimenti per lei (è un amore giovanile ed estivo)
\end{document}
